\newcommand{\latticeToMLB}{%
\begin{figure}[t]
\centering
\resizebox{0.85\textwidth}{!}{%
\begin{tikzpicture}[
  >=Latex,
  line cap=round, line join=round,
  thick,
]
  % --- styles
  \colorlet{nodegreen}{green!60!black}
  \tikzset{
    gedge/.style={draw=nodegreen, line width=1.8pt},
    gnode/.style={draw=nodegreen, fill=white, line width=1.8pt},
    arrow/.style={->, line width=1.8pt, draw=black},
    lbl/.style={font=\Large},
  }

  % --- coordinates (tree on the left)
  \coordinate (n1) at (0,2.8);
  \coordinate (n2) at (-1.6,1.4);
  \coordinate (n3) at ( 1.6,1.4);
  \coordinate (n4) at ( 0.8,0.0);
  \coordinate (n5) at ( 2.4,0.0);

  % --- green edges
  \draw[gedge] (n1) -- (n2);
  \draw[gedge] (n1) -- (n3);
  \draw[gedge] (n3) -- (n4);
  \draw[gedge] (n3) -- (n5);

  % --- green open nodes
  \draw[gnode] (n1) circle [radius=0.30];
  \draw[gnode] (n2) circle [radius=0.28];
  \draw[gnode] (n3) circle [radius=0.28];
  \draw[gnode] (n4) circle [radius=0.26];
  \draw[gnode] (n5) circle [radius=0.26];

  % --- labels (node numbers)
  \node[lbl, anchor=south] at ($(n1)+(0,-0.85)$) {$p$};
  \node[lbl, anchor=north] at ($(n2)+(0,-0.25)$) {$x$};
  \node[lbl, anchor=north] at ($(n3)+(0,-0.25)$) {$q$};
  \node[lbl, anchor=north] at ($(n4)+(0,-0.25)$) {$y$};
  \node[lbl, anchor=north] at ($(n5)+(0,-0.25)$) {$z$};

  % --- middle arrow + label
  \coordinate (Astart) at (3.6,1.4);
  \coordinate (Aend)   at (6.6,1.4);
  \draw[arrow] (Astart) -- (Aend);
  \node[font=\bfseries\Large] at ($(Astart)!0.5!(Aend)+(0,0.6)$) {\ac{mlb}};

  % --- right-side explanation (siblings + uncle)
  \node[anchor=west, font=\Large] at (7.1,1.8) {$(y,z,x)$};
  \node[anchor=west, font=\Large] at (7.1,1.1)
    {\small not $(y,z,q)$ because $q$ is the parent/meet of $y$ and $z$;};
  \node[anchor=west, font=\Large] at (7.1,0.5)
    {\small not $(x,q,p)$ because $p$ is the parent/meet of $x$ and $q$.};

\end{tikzpicture}%
}
\caption{Building \ac{mlb} constraints using the ``siblings + uncle'' rule, where siblings $y,z$ are connected to the same parent node $q$ and $x$ is the uncle node.}
\label{fig:concept_lattice2mlb}
\end{figure}%
}
